\subsection{Skewでない標準CY chartの一般点の除外}
\label{sec:pym-generic-cubic-exclusions}
単一線型条件は，skewでない族でも適用できる．
ここではPym\cite[\S\S3.2--3.6]{Pym}の\emph{指定された標準chart}を扱い，
そのKoszul AS4性・CY性と一般点の表示を既知入力とする．
Pymの六つのgraded CY変形族が全AS4包絡を尽くすとは仮定しない．

\begin{proposition}[四つの非skew CY chartの一般点はH3包絡にならない]
\label{prop:pym-four-generic-empty}
Pymの標準$L(1,1,2)$，$R(2,2)$，$S(2,3)$ chartの非空稠密開部分，
および標準$E(3)$の非零変形は，$Z(B)_3=0$を満たす．
従ってそれらはH3の二次包絡にならない．
特殊パラメータ，chartの境界，成分間の交叉はこの命題では除外しない．
\end{proposition}
\begin{proof}
全てCYなので，系\ref{cor:CY-cubic-central}から$Z(B)_3=0$を示せばよい．

\emph{$L(1,1,2)$．}
Pymの式(5)では$x_0,x_1$は可換，
\[
 x_2x_0=p_0^{-1}x_0x_2,\quad x_3x_0=p_0x_0x_3,\quad
 x_2x_1=p_1x_1x_2,\quad x_3x_1=p_1^{-1}x_1x_3,
\]
かつ$x_3x_2=p_0^{-1}p_1x_2x_3+F(x_0,x_1)$である．
$x_0,x_1$の補助weightを0，$x_2,x_3$を1とすると，PBW性から
associated gradedは
\[
 (q_{01},q_{02},q_{03},q_{12},q_{13},q_{23})
 =(1,p_0,p_0^{-1},p_1^{-1},p_1,p_0p_1^{-1})
\]
のskew環である．$p_0,p_1$が独立な一般パラメータなら，中心単項式の指数$a$は
$a_2=a_3$，$a_0=a_3$，$a_1=a_3$を満たす必要がある．
総次数3では不可能なので，このassociated gradedの中心三次は零である．
中央元の最高weight記号はassociated gradedの中央元となるため，$Z(B)_3=0$となる．
次数3の有限個の指標等式を除けばよいので，これは非空開条件である．

\emph{$R(2,2)$．}
Pymの式(7)のSklyaninパラメータを
$(q_1,q_2,q_3)=(2,3,-5/7)$とする．
$q_1+q_2+q_3+q_1q_2q_3=0$を満たし，標準AS4・CY locus上にある．
六二次関係から生成される関係部分空間のrankを$\mathbb Q$上で求めると，
次数3で44，次数4で221となり，商の次元は20と35である．
この商基底で四つの交換子$Gx_i-x_iG$を並べた140行20列行列はrank20である．
付属証明書に指定した20行の小行列式は
\[
 280000000000/279926361\ne0
\]
となる．従ってこの点の中心三次は零であり，AS-flatな標準族の
中心核の上半連続性から非空開部分でも零となる．
ここでAS正則性はPymの標準Sklyanin族を入力とし，この低次数計算だけから推論しない．

\emph{$S(2,3)$．}
このchartは$B=k[x,y,z][T;\delta]$という次数一のOre拡大であり，
$\delta$はdivergence-freeな二次ベクトル場である．
$\delta\ne0$なら中央元$\sum_{j=0}^m f_jT^j$に対し，$m>0$のとき
基底変数との交換子の最高$T$次数係数は$m f_m\delta(x_i)$となる．
ある$\delta(x_i)\ne0$を選べるので，char0と整域性から矛盾し，中央元は$k[x,y,z]$内にある．
従って$Z(B)_3=\ker(\delta:k[x,y,z]_3\to k[x,y,z]_4)$である．
Pymの式(11)で$b_i=c_i=-1,d_i=0$を選ぶと
\[
 \delta x=x^2-xy-xz,\quad
 \delta y=y^2-yz-xy,\quad
 \delta z=z^2-zx-yz
\]
となり，確かにdivergenceは零である．この15行10列の行列の
指定10行小行列式は$-243$である．よって単射であり，一般点でも単射となる．

\emph{$E(3)$．}
Pymの式(12),(14)のCGG積を，通常の多項式空間上で用いる．
\[
 X=(-5x_0\partial_0-x_1\partial_1+3x_2\partial_2+7x_3\partial_3)/4,
 \qquad Y=4(x_0\partial_1+x_1\partial_2+x_2\partial_3),
\]
\[
 f\star g=\sum_{k\ge0}Y^k(f)\binom Xk(g).
\]
$Y$は各斉次部分で冪零である．$x_0\star f=x_0f$，
$f\star x_0=x_0\sum_k\binom{-5/4}{k}Y^k(f)$なので，
中央元$f$は$Yf=0$を満たす．実際，右の作用素から恒等を引くと
$Y$に定数項$-5/4$の可逆な$Y$多項式を掛けたものとなる．
続いて$x_1\star f=x_1f+4x_0Xf$，$f\star x_1=fx_1$なので$Xf=0$となる．
逆向きもCGG積から従い，中心は$\ker X\cap\ker Y$である．
総次数$n$の単項式の$4X$ weightは$-5n+4\sum_j j a_j$であり，
零weightが存在するには$4\mid n$が必要である．特に$n=3$では中心は零である．
非零変形パラメータはPymの表示により同型なので，全て同じ結論となる．
\end{proof}

\paragraph{検算と分類上の意味．}
\texttt{computations/verify\_cubic\_pym\_generic.py}は，
$R(2,2)$の商空間を自由語から有理数上で構成し，関係の消去と非零小行列式を検証する．
$S(2,3)$の微分行列，$L(1,1,2)$の二つの独立パラメータの整数weight，
$E(3)$の全20個の三次weightも検査する．
これらは有限体の標本ではなく，特性零の証明書である．

Skew $L(1,1,1,1)$の一般点も既に除外されている．
従ってPymの六つの標準CY chartの中で，一般点がH3包絡となり得るものとして
残るのは$R(1,3)$である．その良い中心三次ペンシルの存在は既存の
中心拡大認識・ambient原稿で扱う．
この一般点の結論から，全H3の$R(1,3)$帰属，特殊locusの不在，
または全AS4包絡の被覆を結論してはいけない．
既に構成した量子Fermat型と非CY符号型を保持する必要がある．
